%!TEX root = ../../../main.tex

\glsunset{NAS}
\glsunset{DLNA}
\glsunset{DHCP}
\glsunset{SBOM}
\chapter{Heimnetz-Schutz und Gesetzlicher Rahmen (matr.-Nr. 5783256)}
\section{Schutzkonzepte im Heimnetzwerk}
Eingriffe in die Firmware verschaffen zwar weitreichende Kontrolle über
Smart-Home-Geräte, setzen aber Spezialwissen voraus und kosten meist die
Herstellergarantie. Für die meisten Verbraucher kommt
daher vor allem Anwendungssouveränität in Betracht. Der \gls{SVRV} versteht darunter die Fähigkeit, digitale
Dienste kompetent, sicher und selbstbestimmt zu nutzen, ohne in die Software
der Geräte einzugreifen \cite[S. 12, ff.]{noauthor_digitale_2017}. Im Heimnetz lassen sich die
Datenflüsse unveränderter Geräte mit zwei Ansätzen kontrollieren, der
Netzsegmentierung und der DNS-Filterung. Anschließend werden die Grenzen
dieses Selbstschutzes betrachtet.

\subsection{Netzsegmentierung und logische Isolierung}
Viele Consumer-IoT-Geräte, von IP-Kameras bis zu smarten Steckdosen, sind
unzureichend abgesichert und werden nur kurz mit Updates versorgt. 
Befinden sie sich in einem flachen, unsegmentierten Netz mit
Laptops, Smartphones oder Netzwerkspeichern (\gls{NAS}), genügt die Übernahme eines
einzigen Geräts als Brückenkopf für laterale Angriffe auf das gesamte Heimnetz. 
Das \gls{BSI} empfiehlt deshalb, IoT-Geräte logisch von
den übrigen Systemen zu trennen. In Unternehmensnetzen ist diese
Trennung seit Langem Standard während sie im Heimnetz dagegen noch eine Ausnahme ist \cite{noauthor_smarthome_nodate}.

Dafür gibt es zwei Wege. Das Gastnetz des Routers lässt sich ohne
Fachkenntnisse zuschalten und isoliert die IoT-Geräte vom übrigen Heimnetz.
\glspl{VLAN} erfordern dagegen verwaltbare Switches und
Firewall-Regeln, erlauben aber festzulegen, welche Verbindungen zwischen den
Segmenten zulässig sind, etwa dass Rechner auf IoT-Geräte zugreifen dürfen,
aber nicht umgekehrt \cite{noauthor_smarthome_nodate}. Tabelle~\ref{tab:netzwerkarchitektur}
stellt die Varianten gegenüber.

\begin{table}[htbp]
  \centering
  \caption{Vergleich von Netzwerkarchitekturen im Heimnetz}
  \label{tab:netzwerkarchitektur}
  \footnotesize
  \setlength{\tabcolsep}{4pt}
  \renewcommand{\arraystretch}{1.2}
  \begin{tabularx}{\textwidth}{@{}>{\raggedright\arraybackslash}p{2.6cm}
      >{\raggedright\arraybackslash}p{3.3cm}
      >{\raggedright\arraybackslash}X
      >{\raggedright\arraybackslash}X@{}}
    \hline
    \textbf{Architektur} & \textbf{Einrichtungs\-aufwand}
      & \textbf{Schutz vor lateralen Angriffen}
      & \textbf{Lokale Gerätefunktionen} \\
    \hline
    Flaches Netz (Standard) & Keiner
      & Kein Schutz; ein übernommenes Gerät gefährdet alle anderen
      & Vollständig erhalten \\
    Gastnetz des Routers & Gering; Zuschaltung im Router
      & Hoch; kein Zugriff auf Laptops und NAS
      & Eingeschränkt, da z.\,B. mDNS nicht netzübergreifend funktioniert \\
    VLAN-Segmentierung & Hoch; verwaltbare Switches und Firewall-Regeln
      & Sehr hoch; Filterung nach Protokollen und Adressen
      & Weitgehend erhalten über Routing-Regeln und Reflektoren \\
    \hline
  \end{tabularx}
  \par\smallskip
  \raggedright\scriptsize
  Quelle: Eigene Darstellung auf Basis von \cite{noauthor_smarthome_nodate}.
\end{table}

Die Segmentierung erzeugt allerdings einen Zielkonflikt zwischen Sicherheit
und Komfort: Smart-TVs oder Netzwerklautsprecher, die Medien per \gls{DLNA} oder
AirPlay von einem lokalen \gls{NAS} abspielen, verlieren nach der Isolierung
Funktionen, sofern nicht Routing-Regeln oder Protokoll-Reflektoren nachhelfen \cite{noauthor_smarthome_nodate}.

\subsection{DNS-Filterung und lokale Sinkholes}
Während die Segmentierung vor Angriffen innerhalb des Netzes schützt, richtet
sich die DNS-Filterung gegen den Abfluss von Daten. Viele Geräte senden im
Hintergrund Telemetrie- und Nutzungsdaten an Hersteller oder Werbenetzwerke,
ohne dass sich dies in den Einstellungen abschalten lässt \cite{noauthor_overview_nodate, noauthor_uberblick_nodate}. 
Ansatzpunkt ist das \gls{DNS}, das Domainnamen in IP-Adressen übersetzt. Fast jede Verbindung eines Geräts
beginnt mit einer solchen Anfrage. Ein \gls{DNS}-Sinkhole wird, meist über die \gls{DHCP}-Einstellungen des
Routers, als zentraler \gls{DNS}-Server des Heimnetzes eingetragen und prüft jede
Namensauflösung gegen Sperrlisten. Legitime Anfragen leitet es an einen
externen \gls{DNS}-Server weiter. Für bekannte Tracking- oder Telemetriedomains
antwortet es dagegen mit einer leeren Adresse, sodass das Gerät keine
Verbindung aufbauen kann und die Daten das Heimnetz nicht verlassen \cite{noauthor_overview_nodate}.

Verbreitet sind zwei Open-Source-Lösungen. Pi-hole läuft etwa auf einem
Raspberry Pi, filtert für alle Geräte im Netz ohne Installation auf den
Endgeräten und zeigt in einem Dashboard, welche Geräte welche Domains anfragen
\cite{noauthor_overview_nodate}. AdGuard Home bietet dieselbe Filterung und unterstützt
zusätzlich verschlüsselte Protokolle wie \gls{DoH} und
\gls{DoT}, sodass Internetanbieter oder Dritte die Anfragen nicht
mitlesen können \cite{noauthor_uberblick_nodate}. An ihre Grenze stoßen beide Werkzeuge bei
Geräten, die fest einprogrammierte IP-Adressen oder eigene \gls{DoH}- beziehungsweise \gls{DoT}-Server nutzen.
In diesen Fällen sind zusätzliche Firewall-Regeln nötig \cite{noauthor_uberblick_nodate}.

\subsection{Grenzen des technischen Selbstschutzes}
Segmentierung und \gls{DNS}-Filterung verschaffen technisch versierten Nutzern ein
hohes Maß an Anwendungssouveränität. Betrieb und Pflege eines Sinkholes, das
Beheben fälschlich gesperrter Domains (\emph{False Positives}) und die
Einrichtung von \glspl{VLAN} erfordern jedoch IT-Kenntnisse und dauerhaften Aufwand.
Der \gls{SVRV} betont, dass verbraucherbezogene Digitale Souveränität nicht
voraussetzen darf, dass alle Bürger zu IT-Fachleuten werden \cite[S. 24, ff.]{noauthor_digitale_2017}. Er stützt sie auf die Grundwerte Wahlfreiheit,
Selbstbestimmung, Selbstkontrolle und Sicherheit und unterscheidet drei Handlungsfelder \cite[S. 18]{noauthor_digitale_2017}:
\begin{itemize}
  \item \textbf{Technologie:} nutzerfreundliche, datensparsame und sichere
        Systeme (\emph{Privacy and Security by Design}),
  \item \textbf{Digitale Kompetenz:} Befähigung der Verbraucher zu
        informierten Entscheidungen und kompetenter Nutzung,
  \item \textbf{Regulierung:} staatlicher Ordnungsrahmen zum Schutz vor
        Marktversagen.
\end{itemize}

Wird die Verantwortung allein den Verbrauchern übertragen, entsteht eine
digitale Zweiklassengesellschaft, in der nur IT-affine Haushalte ihr Netz
wirksam abschirmen \cite[S. 31]{noauthor_digitale_2017}. Der SVRV hebt in diesem Zusammenhang hervor, 
dass echte digitale Souveränität nicht durch nachträglichen Konfigurationsaufwand der Nutzer erzwungen werden darf, 
sondern zwingend herstellerseitige Standardeinstellungen voraussetzt, 
die auch technisch nicht versierte Verbraucher im Alltag automatisch schützen \cite[S. 18]{noauthor_digitale_2017}.
Technischer Selbstschutz bleibt damit eine Notlösung, die fehlende Produktsicherheit nicht ersetzt. 
Flächendeckend lässt sich Digitale Souveränität nur durch Regulierung herstellen.

\section{Gesetzlicher Rahmen und herstellerseitige Regulierung}
Beim Kauf können Verbraucher Sicherheit und Datenschutz eines
Smart-Home-Geräts kaum beurteilen. Diese ungleiche Informationsverteilung nimmt
Herstellern den Anreiz, in sichere Architekturen und lange Softwarepflege zu
investieren \cite{noauthor_cyber_nodate,noauthor_regulation_2024}. Die Europäische Union begegnet dem mit dem
\gls{CRA} und das \gls{BSI} ergänzend mit dem
IT"=Sicherheitskennzeichen.

\subsection{Der Cyber Resilience Act der Europäischen Union}
Die Verordnung (EU) 2024/2847 schafft erstmals EU-weit verbindliche
Cybersicherheitsanforderungen für Produkte mit digitalen Elementen. 
Für Elektrogeräte galten bislang zwar strenge Vorgaben etwa zum
Brandschutz, für ihre Cybersicherheit aber kaum Mindeststandards. Der \gls{CRA}
knüpft an die CE-Kennzeichnung an und ist seit Dezember 2024 in Kraft. Die
Meldepflichten gelten seit September 2026 und verpflichten Hersteller, 
aktiv ausgenutzte Schwachstellen binnen 24 Stunden über die neue zentrale Meldeplattform der ENISA an die nationalen CSIRTs zu berichten.
Die übrigen Pflichten gelten ab Dezember 2027.
Hersteller, Importeure und Händler müssen die Anforderungen über den gesamten Lebenszyklus eines Produkts erfüllen \cite{noauthor_regulation_2024}.
Tabelle~\ref{tab:cra_anforderungen} fasst die wichtigsten Pflichten zusammen.

\begin{table}[htbp]
  \centering
  \caption{Kernanforderungen des Cyber Resilience Act und ihr Nutzen für Verbraucher}
  \label{tab:cra_anforderungen}
  \footnotesize
  \setlength{\tabcolsep}{4pt}
  \renewcommand{\arraystretch}{1.2}
  \begin{tabularx}{\textwidth}{@{}>{\raggedright\arraybackslash}p{2.8cm}
      >{\raggedright\arraybackslash}X
      >{\raggedright\arraybackslash}X@{}}
    \hline
    \textbf{Element} & \textbf{Pflicht des Herstellers}
      & \textbf{Nutzen für die Digitale Souveränität} \\
    \hline
    Security by Default
      & Keine voreingestellten Standardpasswörter, automatische
        Sicherheitsupdates \cite{noauthor_regulation_2024}
      & Schutz neuer Geräte ohne eigenen Einrichtungsaufwand \\
    Supportzeitraum
      & Sicherheitsupdates für die erwartete Nutzungsdauer, in der Regel
        mindestens fünf Jahre \cite{noauthor_regulation_2024}
      & Längere sichere Nutzung, weniger softwarebedingte Obsoleszenz \\
    Software-Stückliste \gls{SBOM}
      & Dokumentation aller verwendeten Softwarekomponenten \cite{noauthor_bsi_nodate}
      & Schnelleres Schließen von Lücken in Bibliotheken \\
    Meldepflicht
      & Meldung aktiv ausgenutzter Schwachstellen und schwerer Vorfälle an
        ENISA und nationale CSIRTs \cite{noauthor_regulation_2024}
      & Frühe Warnung und koordinierte Abwehr \\
    \hline
  \end{tabularx}
  \par\smallskip
  \raggedright\scriptsize
  Quelle: Eigene Darstellung auf Basis von \cite{noauthor_regulation_2024} und \cite{noauthor_cyber_nodate}.
\end{table}

Security by Default heißt konkret, dass Geräte ohne Standardpasswörter
ausgeliefert werden und Sicherheitsupdates automatisch einspielen, wobei
Nutzer dies abwählen können \cite{noauthor_regulation_2024}. Die Software-Stückliste \gls{SBOM},
deren Format das \gls{BSI} in der Technischen Richtlinie TR-03183 konkretisiert,
verbleibt beim Hersteller. Sie erlaubt ihm aber, bei neuen Schwachstellen in
Bibliotheken betroffene Produkte schnell zu erkennen \cite{noauthor_bsi_nodate}.
Den Nachweis der Konformität erbringen Hersteller der meisten
Consumer-Geräte durch Selbstbewertung. Dabei ist eine umfassende Risikoanalyse sowie technische Dokumentation erforderlich
und Marküberwachungsbehörden wie die \gls{BSI} sind befugt, dies einzufordern und zu überprüfen.
Für sicherheitskritische Produktklassen sind strengere Verfahren wie die Einbindung von notifizierten Drittstellen vorgesehen \cite{noauthor_regulation_2024}.

Besonders der Supportzeitraum knüpft an die in
Kapitel~\ref{chap:cloud-folgen} beschriebene softwarebedingte Obsoleszenz
an. Hersteller müssen ihn an der zu erwartenden Nutzungsdauer ausrichten und
sein Ende beim Kauf deutlich angeben \cite{noauthor_regulation_2024}. Ein Gerät, das
schon nach zwei Jahren keine Sicherheitsupdates mehr erhält, ist damit nicht
mehr zulässig. Hinsichtlich der Dauer stellt der CRA zudem klar, dass die Regelfrist von fünf Jahren kein starres Maximum bildet, 
sondern bei langlebigen Smart-Home-Komponenten, wie Routern oder zentralen Gateways, an deren tatsächlich zu erwartende Nutzungsdauer angepasst werden muss und somit auch deutlich länger ausfallen kann \cite{noauthor_regulation_2024}.
Der \gls{CRA} regelt somit die technische Sicherheit, garantiert aber nicht die Funktionsfähigkeit von Cloud-Diensten.
Ein Smart-Home-Gerät kann unter dem \gls{CRA} juristisch vollkommen sicher sein, nützt dem Verbraucher aber nichts, 
wenn der Hersteller die Server abschaltet und das Gerät ohne Cloud-Verbindung den Dienst verweigert. 
Das ökonomische Problem des Vendor-Lock-ins und des Funktionsverlusts bei Cloud-Abschaltungen wie es bei Gigaset der Fall war,
wird durch den \gls{CRA} nicht gelöst.

\subsection{Das BSI-IT-Sicherheitskennzeichen}
Während der \gls{CRA} Mindeststandards für das Inverkehrbringen setzt, ist das
IT"=Sicherheitskennzeichen des \gls{BSI} ein freiwilliges Transparenzinstrument am
Point of Sale. Hersteller beantragen es mit der Erklärung, dass
ihr Produkt die Anforderungen des \gls{BSI} erfüllt. Für Smart-Home-Geräte stützen
sich diese vor allem auf die Norm ETSI EN~303~645. Diese fokussieren
sich auf grundlegende Sicherheitsstandards, um dem aktuellen Stand der Technik gerecht zu werden. 
Ein QR-Code auf der Verpackung führt zu einer herstellerunabhängigen Seite des \gls{BSI}. Sie
gibt an, bis wann der Hersteller Sicherheitsupdates zusagt und ob dem \gls{BSI}
ungeschlossene Sicherheitslücken bekannt sind. Das Kennzeichen
ergänzt den \gls{CRA}. Dieser verbietet ab 2027 unsichere Produkte, das Kennzeichen
belohnt schon heute Hersteller, die freiwillig höhere Standards und längere
Updatezeiträume zusagen \cite{noauthor_it-sicherheitskennzeichen_nodate}.

Zusammen entlasten \gls{CRA} und IT"=Sicherheitskennzeichen die Verbraucher von der
Pflicht zum ständigen Selbstschutz. Sie ermöglichen informierte
Kaufentscheidungen und machen grundlegenden Schutz im Betrieb zur Regel statt
zur Ausnahme.
